% A mathematical proof component. Requires amsmath, amssymb, amsthm,
% booktabs and url (or hyperref). This file has no private macro dependency.
\section{The iPRA proof for explicit-binder AX11GH and AX11GHS}
\label{sec:new-relative-ipra}

This section supplies the new proof obligation.  The common numerical
calculus and its calibration to Peano arithmetic are inherited from the
\emph{latest incoming} archive
\path{pa194_ax11h_research_2026-11-14.zip}.  That archive contains a written
weak-arithmetic proof for PA194QAX11H, including the physical and source
bridges.  Its SHA-256 is
\begin{center}\small
\nolinkurl{93ee082ef035b2d32e71d6af5239b25b4bea1773f03bea4d2c39810a7fe15075}.
\end{center}
The older placed archive with the same basename is not the dependency.
The present proof establishes an iPRA halting equivalence with the
calibrated QAX11H table.  Consequently it uses the inherited PA compiler
and the intuitionistic transfer theorem once, without introducing a new
PA presentation or a new reflection principle.

Throughout, halting means halting from the all-zero, two-way infinite
binary tape with the declared start state and head position.  The halt
control is separate from the 194 working controls.  All runs are extended
by an absorbing halt configuration.  The consistency sentence is the
inherited \path{nativePAConsistency}, the negation of the mapped native
Foundation PA-refutation sentence.  The two endpoint files defining this
binding are unchanged between the baseline's ProveIt revision
\path{1c9524d8bf2e0a27b9a16ff622359927f536e663} and the corpus revision
\path{1e1a41f799d4658b829cdd0eabae3775a0c8c4a0}.

\subsection{The common calculus and the changed command}

Put
\[
 p(a,b)=\frac{(a+b)(a+b+1)}2+b,\qquad
 I(a,b)=4p(a,b),\quad M(a,b)=I(a,b)+1,
\]
\[
 U(a,B)=I(B,a)+2,\qquad E(a,b)=I(a,b)+3.
\]
The designated truth and contradiction codes are 0 and 1.  Write
$a\to b=I(a,b)$, $\neg a=I(a,1)$,
$a\wedge b=\neg I(a,\neg b)$, and
$\exists_a B=\neg U(a,\neg B)$.  The numerical calculus
$\mathcal C_4^Q$ has modus ponens, theorem generalization
$B\mapsto U(i,B)$, and the following axiom families:
\begin{align*}
 h_1(D,H)&=(\exists_D U(D,H))\to H,\\
 h_2(D,A,H)&=U(D,U(A,H))\to U(A,U(D,H)),\\
 h_3(D,A,H)&=M(H,D)\to(E(H,A)\to M(A,D)),\\
 h_4(D,A,H)&=E(H,D)\to(M(A,H)\to M(A,D)),\\
 h_5(D,H)&=\exists_D E(H,D),\\
 h_6^Q(D,A,H)&=U(D,I(A,H))\to(U(D,A)\to U(D,H)),\\
 \Gamma(B)&=I(B,U(B,B)),\\
 \rho_8(D,A,H,u)&=((H\to u)\to D)\to((D\to H)\to(A\to H)),
 \quad u\le D.
\end{align*}
Its remaining guarded gate is
\begin{align*}
 V(a,b,c)&=M(a,b)\to U(c,M(a,b)),\\
 P(a,b,c)&=c\to M(c,b),\\
 S(a,b,c)&=U\bigl(b,(M(a,b)\to P(a,b,c))\wedge(M(c,b)\to c)\bigr),\\
 Q(a,b,c)&=U(c,S(a,b,c)\to\neg P(a,b,c)),\\
 G(a,b,c)&=\neg(V(a,b,c)\to Q(a,b,c)),\qquad c\ne a,b.
\end{align*}
The fixed instance $\rho_8(0,0,0,0)=0$ proves the zero code.

Let $[a;w]=p(a,w)$; brackets denote nested Cantor pairs, not free lists.
In particular $[0;0]=0$ and $\operatorname{hd}(0)=\operatorname{tl}(0)=0$.
The store is $(t,w,d,h)$.  Both new machines implement the same commands.
For $A=c$, $D=d+x$ and $H=h+y$, define
\[
 \mathsf A_{c,x,y,z}(t,w,d,h)
       =(0,[q_0,q_1,\ldots,q_9,t;w],z,0),                 \tag{G-A}
\]
where
\[
\begin{array}{c|l@{\qquad}c|l}
 j&q_j&j&q_j\\\hline
 0&G(H,D,A)\text{ if }A\ne H,D;\text{ else }0
     &1&\Gamma(A)\\
 2&h_6^Q(D,A,H)&3&h_5(D,H)\\
 4&h_4(D,A,H)&5&h_3(D,A,H)\\
 6&h_2(D,A,H)&7&h_1(D,H)\\
 8&\rho_8(D,A,H,d)&9&U(c,h).
\end{array}
\]
The sole abstract difference from QAX11H is $q_9=U(c,h)$ in place
of $U(d,h)$.  The binder comes from the explicit AX operand, while
the body still comes from the incoming cache.  The fourth operand is
real: the AX fallthrough consumes a full MP parameter and leaves offset
$z$, even for $z>0$.

The other commands are
\[
 \mathsf S(t,w,d,h)=
 \begin{cases}
  \mathsf{halt},&h=1,\\
  (h,\operatorname{tl}w,d,\operatorname{hd}w),&h\ne1,
 \end{cases}                                             \tag{G-S}
\]
and, with $b=d+x$,
\[
 \mathsf{MP}_x(t,w,d,h)=
 \begin{cases}
  (b,w,b,0),&I(t,b)=h,\\
  (0,w,b,0),&I(t,b)\ne h.
 \end{cases}                                             \tag{G-M}
\]
Only Shift halts.  Every generated $q_0,\ldots,q_8$ is divisible by
four, and $q_9$ has residue two and is at least two.  Thus none is one.

\subsection{Explicit compiler and its weak-base justification}

Write $A_0=\mathsf A_{0,0,0,0}$ and execute words left to right.
The total projection equations and the preceding residue facts give
\begin{align*}
 \mathsf A_{c,x,y,0};\mathsf S^{10}(t,w,d,h)
    &=(q_8,[t;w],0,q_9),\\
 \mathsf A_{c,x,y,0};\mathsf S^{11}(t,w,d,h)
    &=(q_9,w,0,t).
\end{align*}
These words never test the saved $t$.  Define
\[
 \mathsf{Push}=A_0;\mathsf S^{10};\mathsf{MP}_0,
 \qquad \mathsf{Pop}=A_0;\mathsf S^{12},
 \qquad \mathsf{Discard}=\mathsf{MP}_0.
\]
Then
\begin{align*}
 \mathsf{Push}(t,w,d,h)&=(0,[t;w],0,0),\\
 \mathsf{Pop}(t,[a;w],d,h)&=
   \begin{cases}\mathsf{halt},&t=1,\\(t,w,0,a),&t\ne1,\end{cases}\\
 \mathsf{Discard}(t,w,0,h)&=(0,w,0,0),\\
 \mathsf{MP}_i(0,w,0,0)&=(0,w,i,0).
\end{align*}
Push's final comparison fails because its operands have residues zero
and two.  Pop's twelfth Shift sees exactly the saved top.  These statements
include the zero-stack alias.

For $0\le j\le8$, the word
\[
 \mathsf{Discard};\mathsf A_{A,D,H,0};
                \mathsf S^{j+2};\mathsf{Pop}^{9-j}          \tag{G-Sel}
\]
sends $(t,w,0,0)$ to $(q_j,w,0,0)$.  The extra Pops preserve
the noncontradictory selected head and remove the remaining generated
heads followed by a saved zero.  For $u\le D$, the word
\[
 \mathsf{Discard};\mathsf{MP}_u;
 \mathsf A_{A,D\mathbin{\dot-}u,H,0};
                       \mathsf S^{10};\mathsf{Pop}          \tag{G-Rho}
\]
selects $\rho_8(D,A,H,u)$ with the same normalized output shape.

The improved generalization word is
\[
 \mathsf{Gen}_i=\mathsf{Push};\mathsf S;
                       \mathsf A_{i,0,0,0};\mathsf S^{11}. \tag{G-Gen}
\]
Its significant boundaries, starting at $(a,w,0,0)$, are
\[
 (a,w,0,0)\longmapsto(0,[a;w],0,0)
 \longmapsto(0,w,0,a)
 \longmapsto\cdots\longmapsto(U(i,a),w,0,0).
\]
The middle Shift tests cache zero.  The following AX consumes $a$ and
clears the cache before another Shift tests it.  Thus the identity holds
also at $a=1$.  The old offset-setting command is unnecessary because
the binder is now $i$ itself.

For a proof tree $\pi$, let $W_\pi$ be its command word.  Axiom leaves
use (G-Sel) or (G-Rho), and a generalization node appends (G-Gen).
At an MP node with major proof $\pi_1$ of $I(a,b)$ and minor proof
$\pi_0$ of $a$, put
\[
 W_\pi=\mathsf{Discard};\mathsf{Push};W_{\pi_1};
 \mathsf{Push};W_{\pi_0};\mathsf{Push};
           \mathsf S;\mathsf S;\mathsf{MP}_b;\mathsf{Pop}.
\]
The initial Discard and Push reserve a zero below both premises.  Before
the final suffix the store is $(0,[a,I(a,b),0;w],0,0)$ unless it
has already halted.  Two Shifts either halt at $a=1$ or expose
the premises; MP returns $b$; Pop removes the reserved zero, clears
the offset and either halts at $b=1$ or returns $(b,w,0,0)$.
Therefore the exact compiler invariant is
\[
 W_\pi(t,w,0,0)=\mathsf{halt}
 \quad\text{or}\quad
 W_\pi(t,w,0,0)=(\operatorname{concl}(\pi),w,0,0).       \tag{G-Cmp}
\]
If $\pi$ refutes the calculus, the word
\[
                 V_\pi=A_0;W_\pi;\mathsf{Pop}             \tag{G-Ref}
\]
halts from \emph{every} incoming store.  The first AX normalizes the
offset and cache; the final Pop detects top one.  A final plain Shift
would instead test zero and would not establish this assertion.

The word lengths are $12$ for Push, $13$ for Pop, $25$ for Gen,
$121-12j$ for (G-Sel), $26$ for (G-Rho), and $53$ of overhead at
an MP node.  If $n$ is the number of nodes in the expanded derivation
tree, induction gives $|W_\pi|\le121n$ and $|V_\pi|\le121n+14$.
Here $n$ is \emph{not} the number of lines of an unexpanded proof DAG.

The universally quantified base stack in (G-Cmp) must not be hidden
inside an impermissible induction.  Use the baseline's explicit bounded
argument.  For a valid $N$-line proof, with $N>0$, compile words bottom-up and let
$L$ and $K$ be the largest word length and command code.  For an absorbing
primitive recursive command transition $\operatorname{next}$, set
\[
 \Phi_K(s)=1+\max\left(s,
             \max_{c\le K,\ u\le s}\operatorname{next}(c,u)\right),
 \qquad B(j)=\Phi_K^{L(N-1-j)}(s_0).
\]
On malformed proof inputs the compiler returns a fixed default word;
thus the compiler and its finite maxima are total primitive recursive
functions.  For $0\le i<j<N$,
$\Phi_K^L(B(j))\le B(i)$.  Induct on the bounded Boolean statement
that every earlier word is correct on every normalized store at most
its assigned bound, retaining all earlier lines rather than only the
immediate predecessor.  A premise call inside word $j$ is reached after
at most $L$ steps.  If execution has already halted, absorption proves
the desired alternative.  Otherwise the preceding local contracts make
the premise input normalized, and its code lies in the bound assigned to
that earlier line.  The local word identities prove each extension.  This is
quantifier-free induction on a primitive recursive characteristic
function.  At the last line $B(N-1)=s_0$, so the arbitrary initial
store is covered.  This proves (G-Cmp) and (G-Ref) in iPRA.

\subsection{Scheduling and proof extraction}

The scheduler starts at $(n,\ell,t,w,d,h)=(2,0,0,0,0,0)$.
For $\ell=0$ it selects raw seed $n$ and increments $n$; otherwise
it selects $\ell-1$ without incrementing $n$.  Unpair tag zero as
Shift, tag one as MP with one field, and tag $c+2$ as AX with fields
$y,x,z$ in that order.  Let $e(\epsilon)=0$ and
$e(C::v)=1+b_C(e(v))$, where
\[
 b_{\mathsf S}(r)=p(0,r),\quad
 b_{\mathsf{MP}_x}(r)=p(1,p(x,r)),\quad
 b_{\mathsf A_{c,x,y,z}}(r)=p(c+2,p(y,p(x,p(z,r)))).
\]
The projections recover all fields including zeros.  Each positive
continuation strictly decreases.  A fresh seed $m$ therefore drains in
at most $m+1$ command invocations unless it halts earlier.

For a calculus refutation $\pi$, candidate
\[
 N_\pi=p\bigl(2,p(0,p(0,p(0,e(V_\pi))))\bigr)\ge3
\]
executes one additional AX followed by the all-store halting word
$V_\pi$.  The absorbing scheduler has halted within
$K(\pi)=(N_\pi+2)^2$ commands.  This primitive recursive bound
follows by summing the drain bounds for candidates $2,\ldots,N_\pi$:
\[
 \sum_{m=2}^{N_\pi}(m+1)
       =\frac{N_\pi^2+3N_\pi-4}{2}<(N_\pi+2)^2.
\]
The first successful candidate may be earlier.  A bounded induction over
candidate counters establishes arrival or an earlier halt, and the
absorbing flag extends that halt to the displayed bound.  No step
assumes that failed candidates reset the store.

Conversely maintain a finite calculus-proof list with indices proving
the top, cache and every canonical stack head.  Zero is supplied by its
fixed proof, including virtual zero heads of a zero stack.  For $w>0$,
$\operatorname{tl}(w)<w$, so stack decoding is bounded by $w$.
AX's only changed annotation is the Gen step taking the cache proof
of $h$ to a proof of $U(c,h)$.  The restricted head is valid because
$d\le d+x$; the gate is either its guarded axiom or zero.  MP and
Shift preserve the other annotations exactly as in the baseline.
A halting Shift exposes a proof of cache one.  All updates and their
validity checks are primitive recursive, and their cumulative invariant
is a bounded Boolean predicate.  Quantifier-free induction thus gives
\[
 \mathrm{iPRA}\vdash
 \operatorname{Halt}(R_G)\ \longleftrightarrow\
              \exists\pi\,\operatorname{Prf}_{\mathcal C_4^Q}(\pi,1),
                                                               \tag{G-Log}
\]
where $R_G$ is the specified abstract scheduler.

\subsection{Binding the new literal sources and binary tables}

The source registers are $(N,T,W,A,D,H,X,Y,L)$.  Relative to
QAX11H, AX11GH changes only the source operand used in the initial
saved generalizer: the decrement at source address 95 uses $A$
instead of $D$, and the increment at address 100 restores $A$
instead of $D$.  The inherited copy template is uniform in its
selected register, so it gives $U(A,H)$ in the saved stack slot.
Every other AX head, the four-field parser, MP and the cached guard
retain their exact arithmetic equations.

AX11GHS makes one additional source change.  At Shift entry, clear
$T$, duplicate the incoming $H$ into $T,X$ while clearing $H$,
increment $H$, and jump to the common comparison.  The operands
are now $(H,X)=(1,h)$ rather than $(h,1)$; both are cleared by
the comparison and equality occurs exactly at $h=1$.  Both variants
have the same live output as (G-S) and the same retained halted file
\[
 \operatorname{Stop}(n,1,w,0,d,0,0,0,\ell).
\]
Erasing only the data of halted files gives the absorbing abstract halt
of (G-S); this erasure commutes with the source step.  In particular,
no assertion about preserved registers is obtained by silently discarding
a live file.  The GHS preparation cost is
\[
 (2t+1)+(4h+1)+2=2t+4h+4,
\]
whereas AX11GH uses $(2t+1)+(7h+2)+2$.  The common comparison
cost is
\[
 C(a,b)=\begin{cases}
 3a+2,&a=b,\\3b+2,&a<b,\\2a+b+4,&a>b.
 \end{cases}
\]
Use $C(1,h)$ for GHS and $C(h,1)$ for GH.  The unequal branch
then performs the same stack unpair and transfer; the equal branch
performs the same halt instruction.

The delivered source-command checker compares every literal instruction
and outgoing edge with the arithmetic templates, reconstructs the CFG
independently from the source tree, and evaluates the entire 190-token AX
body in a free algebra.  It checks both gate branches, all scratch
registers and the real fourth field.  The template equations and exact
clocks are proved by the baseline's quantifier-free counter inductions;
the two replacements just described instantiate the same copy,
duplicate and comparison lemmas.  Therefore a source command segment
has the claimed endpoint at its specified positive primitive recursive
clock for every coded input, not only the tested examples.

The physical layer uses the baseline's symbolic-stream proof: literal
bit steps, unary-run induction, four list-scan lemmas, exact tape-frame
closure and positive clocks.  Those proofs are parametric in the
source, while the finite binding to each new table is replayed by the
independent literal-table checker.  Its results are
\begin{center}
\begin{tabular}{lrr}
\toprule
 & AX11GH & AX11GHS\\
\midrule
Working controls / transition rows &194 / 388&194 / 388\\
Primitive source sites &3671&3668\\
Certified source outcomes &4778&4774\\
Independent symbolic steps &1106693&1110639\\
Whole blank-start initialization steps &57922&57768\\
Smallest segment-duration constant &3&3\\
\bottomrule
\end{tabular}
\end{center}
Every certificate ranges over nine arbitrary natural primary registers,
an arbitrary nonempty finite spare-register list and an arbitrary
nonnegative gap.  Full tape equality is checked at startup, including
the blank exteriors.  The dense table hashes are
\begin{center}\small
AX11GH:\quad
\nolinkurl{513e9c8b28e0dd41e6003bb8d82c80ff3dac5fb3c5274cd2f709a3bba8c3e9e1}\\[3pt]
AX11GHS:\quad
\nolinkurl{7a4738a9eedd3f53674d67568b3197d38237d11d0fc08fa3d41251a8faf8052c}.
\end{center}

For completeness, the weak-base transport uses the following elementary
principle.  Let $s_k$ be a primitive recursive absorbing source run and
$F$ a primitive recursive frame map such that $F(s)$ is live exactly
when $s$ is live; both halt flags are decidable.  Suppose every live
source step, starting from target configuration $F(s_k)$, has a
specified positive primitive recursive target duration $\delta(s_k)$
and exact endpoint $F(s_{k+1})$, while halted sources have duration
zero and canonical target halt as their frame.  Set
\[
 \tau(0)=0,\qquad \tau(k+1)=\tau(k)+\delta(s_k).
\]
Quantifier-free induction gives
$\operatorname{Run}_{\rm target}(T_0+\tau(k))=F(s_k)$,
where $T_0$ is the verified initialization length.  If $s_k$ is
live, all its predecessors are live by absorption, so $\tau(k)\ge k$.
Given a target halt at time $t$, split on the decidable halt flag of
$s_{t+1}$.  Its live case would give a live certified frame at time
$T_0+\tau(t+1)>t$, contradicting target absorption.  Its halted case
provides the explicit source witness $t+1$.  This is a finite decidable
case split, not double-negation elimination for an unbounded existential.
No minimization, unbounded search or induction with an arbitrary
universal parameter predicate is needed.

Apply this principle first to source-command segments and then to the
physical source-step segments, keeping all frame parameters in the
recursively updated configuration.  The two startup constants must be
kept distinct.  More explicitly, let $z_i$ be the abstract scheduler run,
$\theta(z_i)$ its source-command duration, and
$\Theta(k)=\sum_{i<k}\theta(z_i)$, with zero duration after halt.
Let $\tau(j)$ be the physical cumulative clock of the extended source
run, whose state includes the evolving gap and spare list.  A calculus
refutation $\pi$ supplies the primitive recursive physical witness
\[
             T_0+\tau\bigl(\Theta(K(\pi))\bigr),\qquad
 T_0=\begin{cases}57922&\text{for AX11GH},\\
                  57768&\text{for AX11GHS}.
       \end{cases}
\]
Conversely, two applications of the preceding positive-duration
argument extract an abstract scheduler halt from any physical halt.
Combined with (G-Log), it proves in iPRA, for each new
literal table $M_G$,
\[
 \operatorname{Halt}_{\rm blank}(M_G)
 \longleftrightarrow\operatorname{Halt}(R_G)
 \longleftrightarrow
       \exists\pi\,\operatorname{Prf}_{\mathcal C_4^Q}(\pi,1). \tag{G-TM}
\]

\subsection{Composition with the inherited PA calibration}

Let $M_Q$ be the pinned QAX11H table in the retained baseline packet.
Its command, source and physical proofs give
\[
 \mathrm{iPRA}\vdash
 \operatorname{Halt}_{\rm blank}(M_Q)
 \longleftrightarrow
       \exists\pi\,\operatorname{Prf}_{\mathcal C_4^Q}(\pi,1).
\]
Equation (G-TM) therefore proves the positive halting equivalence of
$M_G$ and $M_Q$ in iPRA.  Intuitionistic logic preserves equivalence
under negation.  The baseline's
\path{proof/consistency_equivalence.tex}, together with its cited
Foundation adapter, finite-demand forward compiler, reverse soundness
and conservation/negative-translation transfer, proves
\[
 \mathrm{iPRA}\vdash
 \operatorname{Con}(\mathrm{PA})
       \longleftrightarrow\neg\operatorname{Halt}_{\rm blank}(M_Q).
\]
Composition yields the claimed theorem for the actual new tables:
\[
 \boxed{\quad
 \mathrm{iPRA}\vdash
 \operatorname{Con}(\mathrm{PA})
       \longleftrightarrow\neg\operatorname{Halt}_{\rm blank}(M_G),
 \qquad M_G\in\{\mathrm{AX11GH},\mathrm{AX11GHS}\}.
 \quad}                                                    \tag{G-Final}
\]
This uses no assumption that PA is consistent and no principle for
extracting a positive existential from a double negation.  It is a
written mathematical proof with explicit primitive recursive algorithms,
local arithmetic invariants and finite implementation certificates.
The full inherited PA calibration is a named, retained proof dependency.
The Python checkers are finite binding audits; they are not a proof
assistant kernel.  No new end-to-end Lean or Rocq verification is claimed.
